\documentclass[11pt]{article}
\usepackage[a4paper,margin=23mm]{geometry}
\usepackage{amsmath,amssymb,xcolor,hyperref,fancyhdr}
\usepackage{fontspec}
\setmainfont{Times New Roman}
\setmonofont{Times New Roman}
\usepackage[italic]{mathastext}
\setlength{\emergencystretch}{3em}
\setlength{\parskip}{7pt}\setlength{\parindent}{0pt}
\pagestyle{fancy}\fancyhf{}\setlength{\headheight}{14pt}
\lhead{FlowMLLab}\rhead{Reading and experiment guide}\cfoot{\thepage}
\begin{document}
\noindent\textbf{Week 8 | Reading the evidence}\par\medskip\hrule\medskip
\section*{Check the branch, the domain and the denominator}
For Rayleigh flow with heat-capacity ratio $\gamma>1$, the supersonic branch has

\[1-\gamma^{-2}<T_0/T_0^*<1.\]

At $\gamma=1.4$ the lower limit is about 0.4898 and corresponds to infinite
Mach number. A requested ratio of 0.30 has no supersonic solution. The solver
must explain this physical limit before invoking a root finder. The sonic endpoint
has $M=1$; no finite maximum Mach should silently redefine the physical domain.

Residual and output error have different units and sensitivities. If a relation
is $F(y,a)=0$, a small perturbation obeys $\delta F\approx F_y\delta y$ with
input $a$ held fixed. A large $F_y$ can produce a large equation residual from a
small output error. A positivity parameterization enforces positivity alone,
not the equation, conservation or accuracy.

Compare baseline error on the common supported queries and report coverage on
the full query set. Error on 65\% of queries is not directly comparable with
error on 100\%. A pass close to a declared threshold is a narrow pass, not a
robust margin. Keep the original criterion when reproducing the experiment.
\paragraph{Check your understanding.}
Identify the input information, the reference, the metric denominator and the
evaluation split in one worked notebook comparison. Predict which conclusion
would change if any one of those four changed, then check your prediction.
\end{document}
